\section{Introdução}

A atividade econômica está associada a externalidades ambientais, entre elas as emissões de dióxido de carbono (CO\textsubscript{2}), que contribuem para as mudanças climáticas. Os efeitos das mudanças climáticas já incluem danos à saúde, aos ecossistemas e à segurança alimentar e hídrica, com agravamento esperado no curto prazo. A redução das emissões integra os esforços globais de mitigação desses efeitos \parencite{ipcc2023spm}.

Este ensaio investiga quais relações intersetoriais compõem as diferenças entre as emissões atribuídas pela produção e pela demanda na economia brasileira de 2015. O objetivo é caracterizar o volume, a concentração e a conectividade dessas relações e discutir suas implicações para a coordenação da mitigação. A análise combina a Matriz de Insumo-Produto (MIP) nacional de 2015, com 67 atividades \parencite{ibge2015mip}, e coeficientes setoriais de emissão, representando as atribuições de emissões como uma rede dirigida e ponderada.

A comparação entre produção e demanda decorre da possibilidade de atribuir as mesmas emissões a diferentes agentes econômicos. No âmbito do Protocolo de Kyoto, a contabilização das emissões nacionais segue uma base territorial, registrando-as no país em que ocorrem \parencite{unfccc2008kyotoManual}. Na atribuição setorial pela produção, as emissões são contabilizadas na atividade que as gera; pela demanda, são associadas ao setor do produto final, considerando os requisitos de sua cadeia produtiva \parencite{marques2012,millerBlair2009}. Por exemplo, as emissões da geração de eletricidade podem ser registradas no setor elétrico ou atribuídas aos produtos finais cuja produção requer essa eletricidade. O critério de atribuição altera a distribuição setorial das emissões, preservando o total contabilizado.

A alocação dos esforços de mitigação depende de quais agentes são responsabilizados por quais emissões. A escolha de um critério de responsabilidade envolve, portanto, uma dimensão normativa, além da identificação das atividades em que as emissões ocorrem. Em termos gerais, questões de justiça exigem reflexão e deliberação pública sobre os valores e os fins que devem orientar a vida em sociedade \parencite{sandel2015justica}. No presente ensaio, a comparação contábil fornece informações para essa discussão, sem determinar como os esforços de mitigação devem ser distribuídos.

A literatura oferece bases para examinar as emissões brasileiras e suas relações com setores, regiões e cadeias produtivas \parencite{imori2015,sanguinetAzzoni2024,montoyaEtAl2026}. No campo das redes, \textcite{huang2019} representam atribuições de emissões na China e combinam filtragem de ligações, indicadores de posição e comunidades. Este ensaio adapta a representação em rede ao recorte nacional brasileiro para identificar quais origens emissoras e quais setores de produtos finais compõem as diferenças entre as contas pela produção e pela demanda.

Os totais e os saldos setoriais não mostram, isoladamente, a composição das relações de atribuição. Um volume elevado de emissões atribuídas pela demanda pode provir de poucas atividades emissoras ou de muitas; um grande emissor pode atender à demanda final de diferentes setores. Além disso, uma diferença pequena entre as duas contas pode resultar de grandes volumes de atribuições em sentidos opostos. A análise de redes permite identificar as contrapartes e distinguir relações concentradas de relações distribuídas entre muitas atividades.

Na rede, cada nó representa uma atividade, e cada ligação dirigida associa uma origem emissora ao setor de um produto final, com peso correspondente ao volume de emissões atribuído. Os volumes e a concentração são examinados antes da filtragem. O filtro de disparidade seleciona ligações destacadas na distribuição local dos pesos, formando uma rede reduzida, denominada \textit{backbone} \parencite{serrano2009}. A concentração dos pesos e a conectividade da rede filtrada descrevem aspectos distintos das atribuições e são interpretadas separadamente.

Os resultados mostram uma rede quase completa, na qual o número de ligações pouco diferencia as atividades, enquanto os pesos revelam diferenças de volume e concentração. A filtragem reduz o número de ligações, mas retém proporções desiguais dos volumes de cada atividade.

O recorte abrange emissões da produção nacional associadas à demanda doméstica e às exportações, excluindo emissões estrangeiras. As intensidades de \textcite{sanguinetAzzoni2024} compõem um cenário exploratório, cuja validação ambiental e monetária permanece pendente, conforme a seção~\ref{sec:dados}. Os resultados não constituem, portanto, um inventário validado das emissões brasileiras de 2015.

A discussão das implicações para a mitigação considera conjuntos de setores e demandas que mereçam investigação conjunta. Intervenções nos processos produtivos e no uso de energia e materiais ao longo das cadeias constituem frentes complementares de mitigação \parencite{rissman2020}. A rede informa quais relações examinar, mas não estima a eficácia nem o custo de instrumentos de política.
